\chapter{Intermediate Representations and Static Single Assignment}
\label{chap:ir}

\section{The NumLang Intermediate Representation Hierarchy}
\label{sec:ir:hierarchy}

To reconcile high-level algebraic metacomputation with native machine code generation, NumLang employs a stratified three-tier intermediate representation pipeline:
\begin{enumerate}
    \item \textbf{High-Level Typed AST}: A structured abstract syntax tree preserving source-level lexical scopes, match expressions, and explicit type annotations (\texttt{src/ast.rs}, \texttt{src/typecheck/typed\_ast.rs}).
    \item \textbf{Control-Flow IR}: A flattened basic-block representation with stack-allocated local variables prior to SSA promotion (\texttt{src/ir/lower.rs}).
    \item \textbf{SSA Mid-Level IR (MIR)}: A strict Static Single Assignment intermediate representation equipped with basic blocks, explicit control-flow terminators, $\phi$-nodes, and MemorySSA tokens (\texttt{src/mir/mod.rs}).
\end{enumerate}

\section{The SSA Mid-Level IR (MIR) Specification}
\label{sec:ir:mir_spec}

The SSA Mid-Level IR is the primary representation upon which the NumLang supercompiler operates. Every function in MIR consists of a set of basic blocks $B_0, B_1, \dots, B_m$, where $B_0$ is the designated entry block.

\begin{definition}[Basic Block Structure]
A Basic Block $B$ consists of:
\begin{equation}
B = \langle \mathtt{id}, \vec{\phi}, \vec{I}, T \rangle
\end{equation}
where $\vec{\phi}$ is an ordered sequence of $\phi$-node assignments, $\vec{I}$ is a sequence of straight-line instructions, and $T$ is a basic-block terminator.
\end{definition}

\subsection{Instruction and Rvalue Grammar}
Instructions compute values or interact with the abstract memory model:
\begin{align*}
\langle Instruction \rangle &::= \mathtt{Assign}(\langle Local \rangle, \langle Rvalue \rangle) \mid \mathtt{Nop} \mid \mathtt{Store}(\langle Local \rangle, \langle Local \rangle) \\
\langle Rvalue \rangle &::= \mathtt{Use}(\langle Local \rangle) \mid \mathtt{Constant}(\langle Int \rangle) \mid \mathtt{BinOp}(\langle Op \rangle, \langle Local \rangle, \langle Local \rangle) \\
&\quad\mid \mathtt{UnaryOp}(\langle UnOp \rangle, \langle Local \rangle) \mid \mathtt{Call}(\langle String \rangle, [\langle Local \rangle]) \\
&\quad\mid \mathtt{Alloc}(\langle Local \rangle) \mid \mathtt{Load}(\langle Local \rangle) \\
&\quad\mid \mathtt{Thunk}(\langle String \rangle, [\langle Local \rangle]) \mid \mathtt{ClosureAlloc}(\langle String \rangle, [\langle Local \rangle])
\end{align*}

\subsection{Control Flow Terminators}
Every basic block in NumLang MIR ends with an explicit terminator that transfers control or returns:
\begin{align*}
\langle Terminator \rangle &::= \mathtt{Return}(\langle Local \rangle) \\
&\quad\mid \mathtt{Branch}(\langle BasicBlockId \rangle) \\
&\quad\mid \mathtt{BranchIf}(\langle Local \rangle, \langle BasicBlockId \rangle, \langle BasicBlockId \rangle) \\
&\quad\mid \mathtt{Switch}(\langle Local \rangle, [(\langle Int \rangle, \langle BasicBlockId \rangle)], \langle BasicBlockId \rangle) \\
&\quad\mid \mathtt{Force}(\langle Local \rangle, \langle Local \rangle, \langle BasicBlockId \rangle) \\
&\quad\mid \mathtt{TypeGuard}(\langle Local \rangle, \langle Type \rangle, \langle BasicBlockId \rangle, \langle BasicBlockId \rangle) \\
&\quad\mid \mathtt{Fork}(\langle BasicBlockId \rangle, \langle BasicBlockId \rangle, \langle BasicBlockId \rangle)
\end{align*}

\section{The MemorySSA Token Model}
\label{sec:ir:memory_ssa}

To enable sound symbolic driving over mutable memory without suffering from global pointer invalidation, NumLang implements a \textbf{MemorySSA} abstraction~\cite{briggs1998effective}.

Every heap operation in MIR consumes and produces an explicit virtual token:
\begin{itemize}
    \item $\mathtt{MemoryDef}(v_{\text{prev}}) \to v_{\text{new}}$: Signifies a store or allocation operation, advancing the memory version.
    \item $\mathtt{MemoryUse}(v)$: Signifies a load or read operation dependent upon memory version $v$.
    \item $\mathtt{MemoryPhi}(v_1, v_2) \to v_3$: Unifies memory versions across converging control-flow paths.
\end{itemize}

During symbolic driving, MemorySSA tokens allow the supercompiler to determine with syntactic precision whether two loads can read the same memory location or whether an intervening store may have overwritten the cell.

\section{Mem2Reg and Alias Analysis}
\label{sec:ir:mem2reg_alias}

\subsection{Briggs SSA Promotion (Mem2Reg)}
Lowering from the high-level AST initially allocates all mutable variables on the stack using explicit \texttt{alloca}, \texttt{load}, and \texttt{store} instructions. The \textbf{Mem2Reg} pass (\texttt{src/mir/mem2reg.rs}) promotes stack allocations to SSA registers:
\begin{enumerate}
    \item \textbf{Dominance Frontiers}: The compiler computes the dominator tree of the CFG and calculates the iterated dominance frontier $DF^+(B)$ for each block containing variable definitions.
    \item \textbf{$\phi$-Node Placement}: For each variable $x$ assigned in block $B$, a trivial $\phi$-node is placed at the entry of every block in $DF^+(B)$.
    \item \textbf{SSA Renaming}: A depth-first traversal of the dominator tree maintains a variable version stack, replacing variable loads with their current dominant SSA register and pruning redundant stores.
\end{enumerate}

\subsection{Field-Sensitive Points-To Analysis}
To support precise heap folding, \texttt{src/mir/alias.rs} executes a flow-sensitive, field-sensitive alias analysis. Pointers are associated with abstract allocation sites:
\begin{equation}
\mathtt{PointsTo}(p) = \{ \langle \mathtt{site}_i, \mathtt{offset}_k \rangle \}
\end{equation}
When the supercompiler drives a load instruction $\mathtt{Load}(p)$, it checks if the points-to set of $p$ is a singleton with an unambiguous allocation site in the current symbolic heap. If so, the load is directly resolved to the symbolic stored term, bypassing heap memory entirely.

\section{Formal Small-Step Operational Semantics of SSA MIR}
\label{sec:ir:small_step}

To provide mathematical grounding for downstream verification, we define the small-step operational transition relation $\step$ over abstract machine configurations:
\begin{equation}
\mathcal{M} = \langle bb, \vec{s}, \mrenv, \mrheap, T \rangle
\end{equation}
where $bb \in \mathtt{BasicBlockId}$, $\vec{s}$ is the remaining instruction list, $\mrenv : \mathtt{Local} \to \mathtt{Val}$ is the local environment, $\mrheap : \mathtt{Addr} \to \mathtt{Val}$ is the heap, and $T$ is the block terminator.

\begin{align}
\text{\bf (Assign-Constant)} \quad & \langle bb, (x = \mathtt{Constant}(c)) :: \vec{s}, \mrenv, \mrheap, T \rangle \step \langle bb, \vec{s}, \mrenv[x \mapsto c], \mrheap, T \rangle \\[0.8em]
\text{\bf (Assign-BinOp)} \quad & \frac{\mrenv[y] = v_1 \quad \mrenv[z] = v_2 \quad v_{\text{res}} = v_1 \;\text{op}\; v_2}{\langle bb, (x = \mathtt{BinOp}(\text{op}, y, z)) :: \vec{s}, \mrenv, \mrheap, T \rangle \step \langle bb, \vec{s}, \mrenv[x \mapsto v_{\text{res}}], \mrheap, T \rangle} \\[0.8em]
\text{\bf (Heap-Alloc)} \quad & \frac{a \notin \text{dom}(\mrheap) \quad \mrenv[y] = v}{\langle bb, (x = \mathtt{Alloc}(y)) :: \vec{s}, \mrenv, \mrheap, T \rangle \step \langle bb, \vec{s}, \mrenv[x \mapsto \mathtt{ptr}(a)], \mrheap[a \mapsto v], T \rangle} \\[0.8em]
\text{\bf (Heap-Load)} \quad & \frac{\mrenv[y] = \mathtt{ptr}(a) \quad \mrheap[a] = v}{\langle bb, (x = \mathtt{Load}(y)) :: \vec{s}, \mrenv, \mrheap, T \rangle \step \langle bb, \vec{s}, \mrenv[x \mapsto v], \mrheap, T \rangle} \\[0.8em]
\text{\bf (Branch-True)} \quad & \frac{\mrenv[c] = \mathtt{true}}{\langle bb, [], \mrenv, \mrheap, \mathtt{BranchIf}(c, bb_1, bb_2) \rangle \step \langle bb_1, \text{resolve\_phi}(bb, bb_1, \mrenv), \mrenv, \mrheap, T_1 \rangle} \\[0.8em]
\text{\bf (Branch-False)} \quad & \frac{\mrenv[c] = \mathtt{false}}{\langle bb, [], \mrenv, \mrheap, \mathtt{BranchIf}(c, bb_1, bb_2) \rangle \step \langle bb_2, \text{resolve\_phi}(bb, bb_2, \mrenv), \mrenv, \mrheap, T_2 \rangle}
\end{align}

\section{Canonical SSA MIR Program Gallery}
\label{sec:ir:gallery}

To observe how source language constructs map to concrete SSA MIR basic blocks, we inspect five canonical lowered functions:

\subsection{1. Arithmetic Stream Summation (\texttt{stream\_sum})}
\begin{lstlisting}[caption={SSA MIR for \texttt{stream\_sum}.}]
fn stream_sum(n: i64) -> i64 {
  bb0:
    let _1: i64 = 0;           // initial sum
    let _2: i64 = 0;           // initial i
    branch bb1;

  bb1:
    let _sum: i64 = phi [bb0: _1], [bb2: _new_sum];
    let _i: i64   = phi [bb0: _2], [bb2: _next_i];
    let _cond: bool = binop lt _i, n;
    branch_if _cond, bb2, bb3;

  bb2:
    let _prod: i64    = binop mul _i, 2;
    let _new_sum: i64 = binop add _sum, _prod;
    let _next_i: i64  = binop add _i, 1;
    branch bb1;

  bb3:
    return _sum;
}
\end{lstlisting}

\subsection{2. Naive List Reverse (\texttt{nrev})}
\begin{lstlisting}[caption={SSA MIR for Recursive \texttt{nrev}.}]
fn nrev(xs: ptr) -> ptr {
  bb0:
    let _tag: i64 = load xs;          // 0 = Nil, 1 = Cons
    let _is_cons: bool = binop eq _tag, 1;
    branch_if _is_cons, bb_cons, bb_nil;

  bb_nil:
    let _nil_res: ptr = alloc 0;
    return _nil_res;

  bb_cons:
    let _head_ptr: ptr = binop add xs, 8;
    let _head: i64 = load _head_ptr;
    let _tail_ptr: ptr = binop add xs, 16;
    let _tail: ptr = load _tail_ptr;
    let _rev_tail: ptr = call nrev, [_tail];
    let _single: ptr = alloc 1;
    store _single, 8, _head;
    let _res: ptr = call append, [_rev_tail, _single];
    return _res;
}
\end{lstlisting}

\subsection{3. Fibonacci Companion Matrix Multiply (\texttt{fib\_matrix})}
\begin{lstlisting}[caption={SSA MIR for \texttt{fib\_matrix}.}]
fn mul_mat(a: ptr, b: ptr) -> ptr {
  bb0:
    let _a00: i64 = load a, 0;  let _a01: i64 = load a, 8;
    let _a10: i64 = load a, 16; let _a11: i64 = load a, 24;
    let _b00: i64 = load b, 0;  let _b01: i64 = load b, 8;
    let _b10: i64 = load b, 16; let _b11: i64 = load b, 24;
    
    let _p00_1: i64 = binop mul _a00, _b00;
    let _p00_2: i64 = binop mul _a01, _b10;
    let _r00: i64   = binop add _p00_1, _p00_2;
    
    let _res: ptr = alloc 32;
    store _res, 0, _r00;
    return _res;
}
\end{lstlisting}

\subsection{4. Double Tree Inversion (\texttt{tree\_flip})}
\begin{lstlisting}[caption={SSA MIR for \texttt{tree\_flip}.}]
fn flip(t: ptr) -> ptr {
  bb0:
    let _tag: i64 = load t, 0;
    let _is_node: bool = binop eq _tag, 1;
    branch_if _is_node, bb_node, bb_leaf;

  bb_leaf:
    let _val: i64 = load t, 8;
    let _leaf: ptr = alloc 16;
    store _leaf, 0, 0;
    store _leaf, 8, _val;
    return _leaf;

  bb_node:
    let _left_ptr: ptr = load t, 8;
    let _right_ptr: ptr = load t, 16;
    let _f_right: ptr = call flip, [_right_ptr];
    let _f_left: ptr  = call flip, [_left_ptr];
    let _node: ptr = alloc 24;
    store _node, 0, 1;
    store _node, 8, _f_right;
    store _node, 16, _f_left;
    return _node;
}
\end{lstlisting}

\subsection{5. Knuth-Morris-Pratt DFA Search (\texttt{kmp})}
\begin{lstlisting}[caption={SSA MIR for \texttt{kmp}.}]
fn kmp_search(pat: ptr, text: ptr) -> bool {
  bb0:
    let _i0: i64 = 0;
    let _j0: i64 = 0;
    branch bb_loop;

  bb_loop:
    let _i: i64 = phi [bb0: _i0], [bb_step: _i_next];
    let _j: i64 = phi [bb0: _j0], [bb_step: _j_next];
    let _loop_cond: bool = binop lt _i, 10;
    branch_if _loop_cond, bb_body, bb_fail;

  bb_body:
    let _p_char: i64 = load pat, _j;
    let _t_char: i64 = load text, _i;
    let _match: bool = binop eq _p_char, _t_char;
    branch_if _match, bb_match, bb_mismatch;

  bb_match:
    let _j_inc: i64 = binop add _j, 1;
    let _is_end: bool = binop eq _j_inc, 3;
    branch_if _is_end, bb_found, bb_step;

  bb_found:
    return true;

  bb_fail:
    return false;
}
\end{lstlisting}

This rich SSA MIR formulation decouples source language syntax from hardware instructions, providing the exact structured control-flow substrate required for the supercompilation passes detailed throughout Part II.

\section{MemorySSA Version Token Algebra}
\label{sec:ir:memory_ssa}

To enable aggressive optimization across heap allocations without introducing unsound pointer aliasing assumptions, NumLang implements \emph{MemorySSA} (Phase 38).

In standard SSA, scalar variables are renamed into single-assignment registers. MemorySSA extends this principle to heap state by tracking abstract memory versions:

\begin{equation}
\mathcal{M} = \langle \mathtt{MemoryVersion}, \mathtt{MemoryDef}, \mathtt{MemoryUse}, \mathtt{MemoryPhi} \rangle
\end{equation}

\begin{itemize}
    \item \textbf{MemoryUse}($v_k$): Attaches to memory reads (such as array loads and struct field dereferences), recording that the read observes memory version $v_k$.
    \item \textbf{MemoryDef}($v_{k+1}, v_k$): Attaches to memory stores, consuming memory version $v_k$ and producing a fresh memory version $v_{k+1}$.
    \item \textbf{MemoryPhi}($v_{\text{res}}, [(bb_1, v_1), \dots, (bb_m, v_m)])$: Merges memory versions at control-flow join points.
\end{itemize}

By explicitly threading memory tokens through SSA def-use chains, the supercompiler proves that memory reads occurring between identical memory versions are purely functional, enabling dead-store elimination and store-to-load forwarding across recursive loops.

\section{The Briggs SSA Mem2Reg Promotion Engine}
\label{sec:ir:mem2reg_impl}

Promoting stack allocations (\texttt{Alloca}) to SSA register values requires computing dominance frontiers and placing $\phi$-nodes. Listing~\ref{lst:mem2reg_impl} details the Briggs SSA promotion algorithm implemented in \texttt{src/mir/mem2reg.rs}:

\begin{lstlisting}[language=Rust, caption={SSA Promotion via Dominance Frontiers and Version Stacks (\texttt{src/mir/mem2reg.rs}).}, label={lst:mem2reg_impl}]

use crate::mir::alias::AliasAnalysis;
use crate::mir::dominance::compute_dominance;
use crate::mir::lower::{MirFunction, Rvalue, Statement};
use crate::mir::{compute_cfg, BasicBlockId, Place, Projection};
use crate::typecheck::types::Type;

/// Computes the Iterated Dominance Frontier (IDF) for a set of defining blocks.
pub fn compute_idf(
    defs: &HashSet<BasicBlockId>,
    dominance_frontiers: &HashMap<BasicBlockId, HashSet<BasicBlockId>>,
) -> HashSet<BasicBlockId> {
    let mut idf = HashSet::new();
    let mut worklist: Vec<BasicBlockId> = defs.iter().cloned().collect();

    while let Some(b) = worklist.pop() {
        if let Some(df_b) = dominance_frontiers.get(&b) {
            for y in df_b {
                if idf.insert(y.clone()) && !defs.contains(y) {
                    worklist.push(y.clone());
                }
            }
        }
    }

    idf
}

/// Identifies candidate memory places suitable for scalar SSA register promotion.
pub fn find_promotion_candidates(func: &MirFunction) -> Vec<Place> {
    let mut non_promotable = HashSet::new();
    let mut all_places = HashSet::new();

    // Any place that uses dynamic array indexing or deref is not a candidate
    for block in &func.blocks {
        for stmt in &block.statements {
            let Statement::Assign(dest, rval) = stmt;
            check_place_promotability(dest, &mut all_places, &mut non_promotable);
            for read_p in collect_all_places(rval) {
                check_place_promotability(&read_p, &mut all_places, &mut non_promotable);
            }
        }
    }

    let mut candidates: Vec<Place> = all_places
        .into_iter()
        .filter(|p| {
            !non_promotable.contains(&p.local)
                && (p.projections.is_empty()
                    || (p.projections.len() == 1
                        && matches!(p.projections[0], Projection::Field(_))))
        })
        .collect();

    candidates.sort_by(|a, b| {
        a.local
            .cmp(&b.local)
            .then_with(|| format!("{:?}", a.projections).cmp(&format!("{:?}", b.projections)))
    });
    candidates
}

fn check_place_promotability(
    p: &Place,
    all_places: &mut HashSet<Place>,
    non_promotable: &mut HashSet<String>,
) {
    for proj in &p.projections {
        match proj {
            Projection::Index(_) | Projection::Deref => {
                non_promotable.insert(p.local.clone());
            }
            Projection::Field(_) | Projection::Payload(_) => {}
        }
    }
    all_places.insert(p.clone());
}

fn collect_all_places(rv: &Rvalue) -> Vec<Place> {
    let mut places = Vec::new();
    match rv {
        Rvalue::Use(p) => places.push(p.clone()),
        Rvalue::BinaryOp(_, p1, p2) => {
            places.push(p1.clone());
            places.push(p2.clone());
        }
        Rvalue::UnaryOp(_, p) => places.push(p.clone()),
        Rvalue::Constant(_) => {}
        Rvalue::Call(_, args) => places.extend(args.iter().cloned()),
        Rvalue::Array(elems) => places.extend(elems.iter().cloned()),
        Rvalue::Struct(_, fields) => {
            for (_, p) in fields {
                places.push(p.clone());
            }
        }
        Rvalue::EnumVariant { fields, .. } => {
            places.extend(fields.iter().cloned());
        }
        Rvalue::Discriminant(p) => {
            places.push(p.clone());
        }
        Rvalue::Phi(incoming) => {
            for (_, p) in incoming {
                places.push(p.clone());
            }
        }
        Rvalue::FnPtr(_) => {}
        Rvalue::ClosureAlloc { captured, .. } => {
            places.extend(captured.iter().cloned());
        }
        Rvalue::Alloc(p) | Rvalue::Load(p) => {
            places.push(p.clone());
        }
        Rvalue::Thunk { env, .. } => {
            for e in env {
                places.push(Place { local: e.clone(), projections: vec![] });
            }
        }
    }
    places
}

/// Executes Mem2Reg SSA promotion: promotes non-aliased stack variables and struct fields
/// to pure SSA registers, inserting block arguments / Phi nodes at iterated dominance frontiers.
pub fn promote_memory_to_registers(func: &mut MirFunction) {
    if func.blocks.is_empty() {
        return;
    }

    let entry = func.blocks[0].id.clone();
    let (preds, succs) = compute_cfg(&func.blocks);
    let all_block_ids: Vec<BasicBlockId> = func.blocks.iter().map(|b| b.id.clone()).collect();
    let dom = compute_dominance(entry.clone(), &preds, &succs, &all_block_ids);

    let candidates = find_promotion_candidates(func);
    if candidates.is_empty() {
        return;
    }

    // Map each candidate place to the set of blocks where it is assigned
    let mut place_defs: HashMap<Place, HashSet<BasicBlockId>> = HashMap::new();
    for block in &func.blocks {
        for stmt in &block.statements {
            let Statement::Assign(dest, _) = stmt;
            if candidates.contains(dest) {
                place_defs.entry(dest.clone()).or_default().insert(block.id.clone());
            }
        }
    }

    // Determine blocks requiring Phi nodes for each candidate
    let mut phi_nodes: HashMap<BasicBlockId, Vec<(Place, Place)>> = HashMap::new(); // block -> vec![(candidate, phi_temp)]
    let mut next_phi_id = 0;

    for candidate in &candidates {
        if let Some(defs) = place_defs.get(candidate) {
            let idf = compute_idf(defs, &dom.dominance_frontiers);
            for b in idf {
                let phi_temp_name = format!("_phi_{}_{}", candidate.local, next_phi_id);
                next_phi_id += 1;
                let phi_place = Place {
                    local: phi_temp_name.clone(),
                    projections: vec![],
                };

                // Add to block arguments
                if let Some(target_block) = func.blocks.iter_mut().find(|blk| blk.id == b) {
                    target_block
                        .arguments
                        .push((phi_temp_name.clone(), Type::I64)); // Type::I64 default scalar
                }

                phi_nodes
                    .entry(b)
                    .or_default()
                    .push((candidate.clone(), phi_place));
            }
        }
    }

    // Build dominator tree children
    let mut dom_children: HashMap<BasicBlockId, Vec<BasicBlockId>> = HashMap::new();
    for b in &all_block_ids {
        if *b == entry {
            continue;
        }
        if let Some(parent) = dom.idom.get(b) {
            dom_children.entry(parent.clone()).or_default().push(b.clone());
        }
    }

    // Dominator tree traversal with renaming stacks
    let mut reaching_values: HashMap<Place, Vec<Place>> = HashMap::new();
    let mut phi_incoming_map: HashMap<(BasicBlockId, Place), Vec<(BasicBlockId, Place)>> =
        HashMap::new();

    let mut stmts_to_remove: HashSet<(BasicBlockId, usize)> = HashSet::new();
    let mut temp_replacements: HashMap<Place, Place> = HashMap::new();

    // Map block positions
    let block_indices: HashMap<BasicBlockId, usize> = func
        .blocks
        .iter()
        .enumerate()
        .map(|(i, b)| (b.id.clone(), i))
        .collect();

    rename_in_dom_tree(
        &entry,
        &dom_children,
        &succs,
        &phi_nodes,
        &candidates,
        func,
        &block_indices,
        &mut reaching_values,
        &mut phi_incoming_map,
        &mut stmts_to_remove,
        &mut temp_replacements,
    );

    // Apply statement removal (stores that have been promoted into SSA values)
    for (bid, b_idx) in &block_indices {
        let block = &mut func.blocks[*b_idx];
        let mut new_stmts = Vec::new();

        // 1. Prepend Phi instructions if any were inserted for this block
        if let Some(phis) = phi_nodes.get(bid) {
            for (candidate, phi_place) in phis {
                let incoming = phi_incoming_map
                    .remove(&(bid.clone(), candidate.clone()))
                    .unwrap_or_default();
                new_stmts.push(Statement::Assign(phi_place.clone(), Rvalue::Phi(incoming)));
            }
        }

        // 2. Retain non-removed statements with rewritten uses
        for (stmt_idx, mut stmt) in block.statements.drain(..).enumerate() {
            if !stmts_to_remove.contains(&(bid.clone(), stmt_idx)) {
                rewrite_statement_places(&mut stmt, &temp_replacements);
                new_stmts.push(stmt);
            }
        }

        block.statements = new_stmts;
    }
}

#[allow(clippy::too_many_arguments)]
fn rename_in_dom_tree(

\end{lstlisting}

\section{Field-Sensitive Points-To Alias Analysis}
\label{sec:ir:alias_impl}

NumLang maintains a field-sensitive alias analysis pass in \texttt{src/mir/alias.rs} to determine memory disambiguation for heap allocations and pointer references:

\begin{lstlisting}[language=Rust, caption={Field-Sensitive Alias Analysis and Points-To Tracking (\texttt{src/mir/alias.rs}).}, label={lst:alias_impl}]

use crate::typecheck::typed_ast::TypedLiteral;
use crate::mir::lower::{MirFunction, Rvalue, Statement};
use crate::mir::{Place, Projection};

#[derive(Debug, Clone, Copy, PartialEq, Eq, Hash)]
pub enum AliasResult {
    /// The two memory places never refer to the same memory location.
    NoAlias,
    /// The two memory places might refer to the same memory location.
    MayAlias,
    /// The two memory places are guaranteed to refer to the exact same memory location.
    MustAlias,
}

impl AliasResult {
    pub fn is_no_alias(self) -> bool {
        self == AliasResult::NoAlias
    }

    pub fn is_must_alias(self) -> bool {
        self == AliasResult::MustAlias
    }

    pub fn is_may_alias(self) -> bool {
        self == AliasResult::MayAlias
    }
}

#[derive(Debug, Clone, Copy, PartialEq, Eq, Hash)]
pub enum ModRefResult {
    NoModRef,
    Ref,
    Mod,
    ModRef,
}

impl ModRefResult {
    pub fn has_ref(self) -> bool {
        matches!(self, ModRefResult::Ref | ModRefResult::ModRef)
    }

    pub fn has_mod(self) -> bool {
        matches!(self, ModRefResult::Mod | ModRefResult::ModRef)
    }
}

#[derive(Debug, Clone)]
pub struct AliasAnalysis {
    /// Constant integer value table for local variables: local name -> constant value
    const_ints: HashMap<String, i64>,
}

impl AliasAnalysis {
    pub fn new(func: &MirFunction) -> Self {
        let mut const_ints = HashMap::new();

        // Scan statements for constant integer definitions
        for block in &func.blocks {
            for stmt in &block.statements {
                let Statement::Assign(dest, rval) = stmt;
                if dest.projections.is_empty() {
                    if let Rvalue::Constant(TypedLiteral::Int(val, _)) = rval {
                        const_ints.insert(dest.local.clone(), *val);
                    }
                }
            }
        }

        AliasAnalysis { const_ints }
    }

    /// Queries the alias relationship between two MIR memory places.
    pub fn alias(&self, p1: &Place, p2: &Place) -> AliasResult {
        // 1. If base locals differ:
        // In NumLang, distinct local variable names represent distinct stack allocations.
        if p1.local != p2.local {
            return AliasResult::NoAlias;
        }

        // 2. Both places have the same root local. Compare projections step-by-step.
        let min_len = p1.projections.len().min(p2.projections.len());
        for i in 0..min_len {
            match (&p1.projections[i], &p2.projections[i]) {
                (Projection::Field(f1), Projection::Field(f2)) => {
                    // Field sensitivity: distinct field names on the same struct are disjoint.
                    if f1 != f2 {
                        return AliasResult::NoAlias;
                    }
                }
                (Projection::Index(idx1), Projection::Index(idx2)) => {
                    // Array constant index analysis:
                    // Check if both indices are known constants
                    let c1 = self.resolve_constant_int(idx1);
                    let c2 = self.resolve_constant_int(idx2);

                    if let (Some(v1), Some(v2)) = (c1, c2) {
                        if v1 != v2 {
                            // Distinct constant array indices are provably disjoint!
                            return AliasResult::NoAlias;
                        }
                    } else if idx1 != idx2 {
                        // Dynamic or unknown indices on the same array base may alias
                        return AliasResult::MayAlias;
                    }
                }
                (Projection::Deref, Projection::Deref) => {
                    // Both dereferenced from the same base
                }
                // Incompatible projection types (e.g. Field vs Index) on the same level
                _ => return AliasResult::NoAlias,
            }
        }

        // 3. All common projections matched identically:
        if p1.projections.len() == p2.projections.len() {
            AliasResult::MustAlias
        } else {
            // One is a sub-path / prefix of the other (e.g., p vs p.x)
            AliasResult::MayAlias
        }
    }

    /// Resolves a place to a known constant integer if possible.
    fn resolve_constant_int(&self, place: &Place) -> Option<i64> {
        if place.projections.is_empty() {
            self.const_ints.get(&place.local).copied()
        } else {
            None
        }
    }

    /// Computes the ModRef effect of a statement with respect to a memory place `target`.
    pub fn modref(&self, stmt: &Statement, target: &Place) -> ModRefResult {
        let Statement::Assign(dest, rval) = stmt;

        let mut modifies = false;
        let mut references = false;

        // Check if destination aliases target
        let dest_alias = self.alias(dest, target);
        if dest_alias != AliasResult::NoAlias {
            modifies = true;
        }

        // Check if any read place aliases target
        let reads = collect_reads_for_alias(rval, dest);
        for read_p in reads {
            let read_alias = self.alias(&read_p, target);
            if read_alias != AliasResult::NoAlias {
                references = true;
                break;
            }
        }

        match (modifies, references) {
            (true, true) => ModRefResult::ModRef,
            (true, false) => ModRefResult::Mod,
            (false, true) => ModRefResult::Ref,
            (false, false) => ModRefResult::NoModRef,
        }
    }
}

fn collect_reads_for_alias(rv: &Rvalue, dest: &Place) -> Vec<Place> {
    let mut reads = Vec::new();

    for proj in &dest.projections {
        if let Projection::Index(idx_place) = proj {
            reads.push((**idx_place).clone());
        }
    }

    match rv {
        Rvalue::Use(p) => reads.push(p.clone()),
        Rvalue::BinaryOp(_, p1, p2) => {
            reads.push(p1.clone());
            reads.push(p2.clone());
        }
        Rvalue::UnaryOp(_, p) => reads.push(p.clone()),

\end{lstlisting}
